\chapter{零点定理的证明}\label{chap:nullstellensatz-proof}

\noindent\textbf{第 \thechapter{} 章学习目标}
\begin{itemize}
  \item 定义环上的模和代数, 并给出例子.
  \item 描述有限生成代数与有限生成模之间的区别.
  \item 利用整性的概念, 判定何时一个有限生成代数实际上是有限生成模.
  \item 通过 Noether 正规化描述有限生成 $K$-代数的一般结构.
  \item 利用 Noether 正规化证明零点定理.
\end{itemize}

我们已经通过代数与几何的等价看到了零点定理的威力, 现在该证明它了. 要走到零点定理的证明, 必须先花相当一段篇幅考察纯代数内容, 包括研究 $R$-模和 $R$-代数, 最后在第~\ref{sec:noether-normalization}~节证明 Noether 正规化定理.

除了是证明零点定理的关键一步外, Noether 正规化定理本身也是一个关于有限生成 $K$-代数结构的重要结论. 它说明, 这样的代数固然不必是有限生成的 $K$-向量空间 (例如多项式环 $K[x]$ 就不是), 却总能表示为有限生成的\textquotedbl{}向量空间\textquotedbl{}, 其标量取自一个同构于多项式环的子代数. 我们给\textquotedbl{}向量空间\textquotedbl{}加上引号, 是因为这里的标量不构成域. 所以, 要理解 Noether 正规化, 就必须推广向量空间的定义, 允许标量来自一般的环 (照例假设它是带单位元的交换环). 这些广义的向量空间称为模, 在第~\ref{sec:r-modules}~节引入.

带单位元的交换环上的模, 其基本定义与域上的向量空间并无不同, 但这个更一般背景下的理论引出了不少新概念, 其中最重要的是整性的概念及其与有限生成的关系. 本章前四节用于发展 $R$-模和 $R$-代数的理论, 最后在第~\ref{sec:noether-normalization}~节给出 Noether 正规化的证明. 在第~\ref{sec:proof-of-nullstellensatz}~节, 这些准备将使我们证明零点定理, 从而完整地论证代数与几何的等价.

\par\smallskip
{\footnotesize
\noindent\textcopyright{} The Author(s) 2025\par
\noindent E. Clader, D. Ross, \emph{Beginning in Algebraic Geometry}, Undergraduate Texts in Mathematics,
\href{https://doi.org/10.1007/978-3-031-88819-9_6}{\nolinkurl{https://doi.org/10.1007/978-3-031-88819-9_6}}.\par}
